% Level 2, Buffer session B1 -- Cipher & Puzzle Day
\kcpart{B1}{Cipher \& Puzzle Day}{Codes hide messages, patterns crack them, and logic unlocks the treasure chest.}
\lessoninfo{Concepts: substitution ciphers (I/R), number codes (I), logic puzzles (I), algorithms (R), dictionaries (R, optional)}{Tools: paper, pencils, scissors; computer optional}{CSTA: 2-NI-05, 2-NI-06, 2-AP-10}

\begin{teachernote}[When to use this session]
Any week. It is fully unplugged, so it works after a missed session, before a holiday, or when screens need a rest.
If your group has not done week 5 yet, this session is a gentle first meeting with the Caesar cipher; if they have, it deepens it
(keyword ciphers and codebreaking).
\end{teachernote}

\begin{objectives}
  \item The student can encode and decode a message with a Caesar shift, a number code and a keyword cipher by following the steps exactly. (L2.G8; L2.WB1.1)
  \item The student can crack a Caesar message without the key by trying shifts in order or by spotting the most common letter. (L2.G2; L2.WB1.2)
  \item The student can solve a logic puzzle by writing down facts and ruling out possibilities one at a time. (L2.G8; L2.WB1.3)
  \item The student can spot a message that tries to trick them into sharing a password or secret. (L2.G9; L2.WB1.4)
\end{objectives}

\begin{materials}
  \item Print: the \textbf{cipher wheel page} from workbook chapter 5 (one per student), and the \textbf{puzzle cards} page at the end of this plan
        (one set per pair; cut along the lines). Keep the answers in this plan with you.
  \item Scissors, split pins or pencils (for the wheels), scrap paper, pencils.
  \item A small \enquote{treasure}: a box or envelope with a 4-digit code written on the lid as \texttt{\_ \_ \_ \_} and a treat or sticker inside
        (optional, but very motivating). The code is \textbf{7253}.
  \item Optional: one Mac with \texttt{grade-4-6/2-intermediate/exercises/week05/ex-caesar-cipher/solution.py} open in Thonny, only if the group has done week 5.
  \item A timer.
\end{materials}

\needspace{16\baselineskip}
\section{At a glance}
Times are a guide: set a timer for each phase and move on when it rings.
\begin{timeplan}
  \timeslot{0--8}{Warm-up}{Build the cipher wheels (or reuse them) and decode a welcome message.}
  \timeslot{8--18}{Concept}{Three kinds of code: shift (Caesar), number (A1Z26), keyword. Why patterns make codes crackable.}
  \timeslot{18--28}{Guided demo}{Crack a Caesar message without the key, step by step, on paper.}
  \timeslot{28--52}{Hands-on build}{Pairs: the four-puzzle escape challenge to open the treasure; bug hunt on a broken decoding; last 6 minutes: invent a puzzle.}
  \timeslot{52--60}{Wrap-up}{Trick-message spotting (phishing), exit ticket, reflection.}
\end{timeplan}

\section{Warm-up (0--8 min)}
\begin{unplugged}[Cipher wheels and a welcome message]
Students cut out the two wheels (or take out the ones they made in week 5) and pin them together. Write on paper:
\texttt{zhofrph frgheuhdnhuv} and say: \enquote{This was made with shift 3. Turn small \textbf{a} under big \textbf{d}. Find each secret letter
on the \textbf{big} wheel and read the letter under it on the \textbf{small} wheel.}
(Answer: \enquote{welcome codebreakers}.)
\end{unplugged}

\section{Concept (8--18 min)}
\begin{teacherscript}[Say]
\enquote{A \textbf{cipher} is a rule for turning a message into secret code. To use one, you follow the rule exactly, step by step, every time:
that is an \textbf{algorithm}, just like a program. Today we use three ciphers.
The \textbf{Caesar shift}: every letter moves the same number of places.
The \textbf{number code}: a is 1, b is 2, all the way to z is 26.
The \textbf{keyword cipher}: you write a secret word first, then the rest of the alphabet, and that is your secret alphabet.
And then we become codebreakers: secret codes have patterns, and patterns give them away.}
\end{teacherscript}
Write the keyword alphabet for the keyword \texttt{DRAGON} on paper while you explain (repeated letters are skipped; the rest follow in order):
\begin{center}
\setlength{\tabcolsep}{2.2pt}\small
\begin{tabular}{r*{26}{c}}
plain: & a&b&c&d&e&f&g&h&i&j&k&l&m&n&o&p&q&r&s&t&u&v&w&x&y&z \\
secret: & \textbf{d}&\textbf{r}&\textbf{a}&\textbf{g}&\textbf{o}&\textbf{n}&b&c&e&f&h&i&j&k&l&m&p&q&s&t&u&v&w&x&y&z \\
\end{tabular}
\end{center}

\begin{analogy}[Footprints in the snow]
A Caesar shift hides each letter, but not the pattern of the letters. In English, \textbf{e} is the most common letter, so in a long secret message
the most common secret letter is probably \textbf{e} in disguise. It is like footprints in snow: you cannot see who walked past,
but you can see where they went. Codebreakers follow the footprints.
\end{analogy}

\section{Guided demo (18--28 min)}
\begin{workedexample}[Cracking a Caesar message without the key]
Write this on paper: \texttt{tlla tl ulhy aol nyllu ayll ha aoyll}
\begin{enumerate}[leftmargin=1.4em]
  \item \textbf{Count letters.} Make a tally together. The most common secret letter is \textbf{l} (11 times).
  \item \textbf{Guess.} \enquote{If \textbf{l} is really \textbf{e}, how far is it from e to l?} Count on the wheel or the alphabet: e, f, g, h, i, j, k, l is \textbf{7} places. So the shift is probably 7.
  \item \textbf{Test on a short word.} \texttt{aol} with shift 7 back: a → t, o → h, l → e: \enquote{the}. A real word, so the guess is good.
  \item \textbf{Decode the rest}: \enquote{meet me near the green tree at three}.
  \item \textbf{Plan B (brute force):} if counting does not work (short messages), try shift 1, 2, 3 \ldots{} on the first word only. At most 25 tries.
\end{enumerate}
Say: \enquote{Both plans are algorithms: exact steps that always work. A computer can do plan B in a blink, which is why real secrets need much stronger codes than Caesar.}
\end{workedexample}
Only if your group has done week 5: open \texttt{week05/ex-caesar-cipher/solution.py} in Thonny, press \menu{Run}, type \texttt{d}, paste the message
and shift 7 to check. (The program is printed in the Answer Key, Exercise 5.2.)

\section{Hands-on build (28--52 min)}
\begin{pairroles}[Codebreaker pairs]
The \textbf{writer} holds the pencil and writes the decoded letters; the \textbf{reader} works the wheel or the key and calls out each letter.
Swap after each puzzle card.
\textbf{1 student:} you are the reader; they write. Swap after each card.
\textbf{2 students:} one pair with one set of cards.
\textbf{3 students:} add a \textbf{checker} who re-checks every third letter; rotate roles after each card.
\textbf{4 students:} two pairs race to the treasure code, each with its own set of cards. Both teams that finish get the treasure!
\end{pairroles}

Hand out the four puzzle cards (end of this plan). Each card gives one digit of the 4-digit treasure code. Help only with questions like
\enquote{What does the card ask you to do?}, not with letters.

\begin{bughunt}[Kai's broken decoding]
Write this on paper: \enquote{Kai decoded \texttt{brx duh d vwdu} with shift 3 and got \texttt{?ou are a star}. Leo decoded \texttt{zhoo grqh} with
shift 3 and got \texttt{ckrr jutk}. What went wrong for each?}
\par\textbf{Answers.} Kai forgot to \textbf{wrap around}: going back 3 from \textbf{b} goes past \textbf{a} to \textbf{y}, so the word is \enquote{you}.
Leo shifted the \textbf{wrong way}: he went 3 forward instead of 3 back (the answer is \enquote{well done}). These are exactly the bugs a
Caesar program can have: missing wrap-around and wrong direction.
\end{bughunt}

\begin{freebuild}[Invent a puzzle]
Last 6 minutes: each pair invents a puzzle card for the other pair (or for someone at home): a secret message in one of today's ciphers,
or a small logic puzzle like card 4. Rule: the solver must be able to solve it with the steps they learned today. Write the answer on the back.
\end{freebuild}

\section{Wrap-up (52--60 min)}
\begin{safety}[Spot the trick message]
Read out these three messages and ask \enquote{Real or trick?}
(1) \enquote{Your game account will be deleted! Reply with your password in 10 minutes to save it.}
(2) \enquote{Mom: dinner is at 6.}
(3) \enquote{You won 1000 free gems! Type your parent's card number to claim them.}
Messages 1 and 3 are tricks (called \textbf{phishing}: fishing for your secrets). Clues: hurry-up pressure, asking for a password or money details,
prizes you did not enter for. Agree: real companies never ask for your password in a message. If in doubt: do not reply, show a grown-up.
\end{safety}

\begin{formativecheck}[Exit ticket]
Decode \texttt{mjqqt} (shift 5). \textbf{Answer:} \texttt{hello}. Bonus: why is a Caesar cipher quick to crack?
(Only 25 shifts to try, and common letters like e give it away.)
\textbf{Look for:} students who shift forward instead of back; revisit the wheel direction.
\end{formativecheck}

\begin{reflection}[Say it or write it]
\begin{itemize}[leftmargin=1.2em]
  \item Which puzzle was hardest? What did you do when you got stuck?
  \item How is following a cipher like running a program?
  \item What makes a secret code strong, and what makes it weak?
\end{itemize}
\end{reflection}

\section{Notes for the teacher}
\begin{commonmistakes}
  \item \textbf{Turning the wheel the wrong way} or reading from the wrong ring. Encode: small → big. Decode: big → small.
  \item \textbf{Forgetting to wrap around} after z (or before a) with a paper alphabet. The wheel avoids this: use it.
  \item \textbf{Keyword alphabet with repeated letters.} Each letter appears only once; skip letters already used in the keyword.
  \item \textbf{Number code: losing word breaks.} Dashes separate letters, spaces separate words.
  \item \textbf{Logic puzzle: guessing.} Encourage a small table with ✓ and ✗ for each fact.
\end{commonmistakes}

\begin{differentiation}
\supportpath Give the Caesar card first and let the student keep the wheel set at the right shift. Write the first decoded word together.
For card 4, draw the table with names down the side and gems along the top.
\extensionpath Use the keyword cipher with a keyword of their own. Then try frequency analysis on a longer message from a book, encoded by a partner.
If the group knows week 5, write \texttt{keyword\_key(word)} in Python that builds the keyword cipher as a dictionary.
\end{differentiation}

\begin{stretchchallenge}[The Vigenere twist]
Instead of one shift for the whole message, use a different shift for each letter, taken from a keyword: with keyword \texttt{abc}
the shifts are 0, 1, 2, 0, 1, 2 \ldots{} Encode \texttt{hello} (answer: \texttt{hfnlp}). Why is this much harder to crack by counting letters?
(The same letter turns into different secret letters.)
\end{stretchchallenge}

\begin{teachernote}[If you are short on time]
Skip building wheels (use the alphabet strip in workbook chapter 5) and the keyword cipher explanation; do cards 1, 2 and 4 only
(the treasure code becomes 7\_3: tell them the third digit is 5). Keep the trick-message talk.
\end{teachernote}

\clearpage
\section{Puzzle cards (cut out, one set per pair)}
\begin{center}
\renewcommand{\arraystretch}{1.4}
\begin{tabular}{|p{0.46\linewidth}|p{0.46\linewidth}|}
\hline
\textbf{Card 1: Caesar}\par
Shift: \textbf{1}. Decode:\par
\texttt{uif gjstu ovncfs jt tfwfo}\par\rule{0pt}{2.2cm}
&
\textbf{Card 2: Number code}\par
a = 1, b = 2, c = 3 \ldots{} z = 26. Decode:\par
\texttt{19-5-3-15-14-4 \ \ 9-19 \ \ 20-23-15}\par\rule{0pt}{2.2cm}
\\ \hline
\textbf{Card 3: Keyword cipher}\par
Keyword: \textbf{DRAGON}. Make the secret alphabet: write DRAGON, then the rest of the alphabet without repeats.
Decode (secret → plain):\par
\texttt{tceqg es nevo}\par\rule{0pt}{1.8cm}
&
\textbf{Card 4: Logic gems}\par
Ada, Leo, Mei and Omar each found one gem: a ruby (red), an emerald (green), a sapphire (blue) and a topaz (yellow).
Ada's gem is not red. Leo found the sapphire. Mei's gem is green.
Who found the topaz? The \textbf{last} digit is the number of letters in that name.\par\rule{0pt}{0.6cm}
\\ \hline
\end{tabular}
\end{center}

\begin{answer}[Puzzle answers (keep this part)]
Card 1: \enquote{the first number is seven} → \textbf{7}.\quad
Card 2: \enquote{second is two} → \textbf{2}.\quad
Card 3: \enquote{third is five} → \textbf{5}.\quad
Card 4: Leo has the sapphire and Mei the emerald; Ada's is not the ruby, so Ada has the topaz and Omar the ruby. \enquote{Ada} has 3 letters → \textbf{3}.
\textbf{Treasure code: 7253.}\par
Warm-up: \enquote{welcome codebreakers}.
\end{answer}
